\section{Related Literature}\label{sec:literature}

This paper contributes to five literatures. First, it speaks to work on representation and party responsiveness by studying a direct, high-frequency channel of supporter input rather than only polls, election returns, manifestos, or roll calls. Second, it builds on political economy research on the internet and media by studying whether a digital political entrepreneur's organization learns from online voice after mobilization occurs. Third, it connects to work on online feedback and agenda setting by shifting attention from politicians on general social media to a party-owned platform during party formation and institutionalization.

The project also contributes to research on M5S, digital parties, and platform politics. Existing work documents the movement's hybrid organization, claims of direct participation, and platform-mediated politics. This paper aims to add a validated longitudinal measure of whether supporter text influenced later agenda and policy. Finally, the paper uses text-as-data and LLM-based measurement as tools, with human validation and explicit uncertainty rather than as substitutes for research design.

Specific citations and literature positioning are TBD. See \texttt{docs/LITERATURE\_MAP.md}.
